\hspace{3mm}

\centerline{\bf \huge 4 класс}
\centerline{\normalsize{\it Работа рассчитана на \bf{90} минут}}

\hspace{3mm}

\problem{4.1} Расставьте скобки так, чтобы получилось верное равенство: 

$2\cdot 74 - 2 : 2 + 16 = 8$

\hspace{3mm}

\problem{4.2} В 9:00 на остановке «Университет» стояли 4 белых автобуса. Через 10 минут между каждым автобусом припарковались по 3 синих автомобиля. Ещё через 15 минут между каждым транспортным средством встали по 2 красных трактора. Сколько стояло транспортных средств на остановке «Университет» в 9:20, если известно, что никакое транспортное средство не уезжало с остановки?

\hspace{3mm}

\problem{4.3} В каждой из двадцати школ города Элисты обучаются поровну учеников. Если в каждой школе будет обучаться на 50 учеников меньше, то во всех школах останется столько учеников, сколько их раньше было суммарно в девятнадцати школах. Сколько учеников изначально обучалось в каждой школе?

\hspace{3mm}

\begin{wrapfigure}{R}{0.3\textwidth}
\vspace{-20mm}
\centering
\includegraphics[width=0.35\textwidth]{Городская Олимпиада Элиста 2023/Условия/picture_4.4.png}
\caption{Рис. 1}
\end{wrapfigure}

\problem{4.4} Однажды Пюрвя вырезал по линиям клетки из клетчатого листа бумаги 4 одинаковые фигурки и решил их склеить друг с другом. У него получилась замечательная фигура (рисунок справа). Определите, как разрезать замечательную фигуру на изначальные 4 одинаковые фигурки Пюрви?

\hspace{3mm}

\problem{4.5} На шахматной доске необычного размера 23 клетки в длину и 23 клетки в ширину, состоящей из 264 белых полей и 265 чёрных, (расположенных в шахматном порядке) стоит по одной белой пешке на каждом белом поле, а на чёрных полях ничего не стоит. В какой-то момент времени Рома решил одновременно взять каждую белую пешку и переставить в какое-то соседнее по стороне поле. Какое минимальное количество чёрных полей будет не занято, если на одном поле может стоять более одной пешки?